\documentclass{notaaula}

\titulonota{Introdução ao magnetismo}
\creditos{3º ano · Física · Setembro/2026 · Prof. Josué Cavalcante}

% --- auxiliares só desta nota: pedaços de ímã e linhas de campo ---
% \pedaco{x esquerdo}{x direito}{y}  — N à esquerda, S à direita
\newcommand{\pedaco}[3]{%
  \pgfmathsetmacro{\pxm}{(#1+#2)/2}%
  \fill[naCoral] (#1,{#3-0.23}) rectangle (\pxm,{#3+0.23});
  \fill[naAzul]  (\pxm,{#3-0.23}) rectangle (#2,{#3+0.23});
  \node[font=\tiny\bfseries, text=naFundo] at ({(#1+\pxm)/2},#3) {N};
  \node[font=\tiny\bfseries, text=naFundo] at ({(\pxm+#2)/2},#3) {S};}
% idem, com os polos trocados
\newcommand{\pedacoI}[3]{%
  \pgfmathsetmacro{\pxm}{(#1+#2)/2}%
  \fill[naAzul]  (#1,{#3-0.23}) rectangle (\pxm,{#3+0.23});
  \fill[naCoral] (\pxm,{#3-0.23}) rectangle (#2,{#3+0.23});
  \node[font=\tiny\bfseries, text=naFundo] at ({(#1+\pxm)/2},#3) {S};
  \node[font=\tiny\bfseries, text=naFundo] at ({(\pxm+#2)/2},#3) {N};}
\tikzset{
  linhaB/.style={naVerde, thick,
    decoration={markings, mark=at position 0.55 with
      {\arrow[naVerde]{Stealth[length=4pt]}}},
    postaction={decorate}},
}

\begin{document}
\cabecalho

\begin{sobre}
Eletromagnetismo · ímãs e polos · campo magnético $\vec{B}$ · linhas de campo ·
bússola e o campo da Terra · 3º A/B/C · FGB.
\end{sobre}

% ============================================================
\section{Ímãs e polos magnéticos}

\begin{copiar}{Noção de ímã}
\definicao \dest{Ímã} é o corpo capaz de atrair pedaços de ferro e de orientar
uma agulha de bússola. Todo ímã tem duas regiões onde o efeito é mais forte,
os \dest{polos}: o polo \dest{norte} (N) e o polo \dest{sul} (S).

\rotulo{Regra da interação.} Entre dois ímãs,
\begin{itens}
\item polos \dest{iguais} (N–N ou S–S) se \dest{repelem};
\item polos \dest{diferentes} (N–S) se \dest{atraem}.
\end{itens}
\diaadia{O ímã de geladeira gruda no aço da porta, mas não no alumínio da
panela nem no cobre do fio.}
\end{copiar}

\begin{figuranota}{Acima, N diante de S: os ímãs se aproximam. Abaixo, N diante
de N: os ímãs se afastam.}
\begin{tikzpicture}
  % atração: S de frente para N
  \pedaco{-3.0}{-1.3}{0.75}
  \pedaco{-0.4}{1.3}{0.75}
  \draw[naVerde, thick, -{Stealth[length=4.5pt]}] (-1.20,0.75) -- (-0.95,0.75);
  \draw[naVerde, thick, {Stealth[length=4.5pt]}-] (-0.75,0.75) -- (-0.50,0.75);
  \node[font=\tiny, color=naVerde, anchor=west] at (1.5,0.75) {atração};
  % repulsão: N de frente para N
  \pedacoI{-3.0}{-1.3}{-0.55}
  \pedaco{-0.4}{1.3}{-0.55}
  \draw[naCoral, thick, {Stealth[length=4.5pt]}-] (-1.20,-0.55) -- (-0.95,-0.55);
  \draw[naCoral, thick, -{Stealth[length=4.5pt]}] (-0.75,-0.55) -- (-0.50,-0.55);
  \node[font=\tiny, color=naCoral, anchor=west] at (1.5,-0.55) {repulsão};
\end{tikzpicture}
\end{figuranota}

% ============================================================
\section{Os polos são inseparáveis}

\begin{copiar}{Não existe monopolo}
\definicao Ao partir um ímã ao meio, \dest{não} se obtém um pedaço só com N e
outro só com S: cada pedaço vira um \dest{ímã completo}, com os dois polos.
Repetindo o corte, o resultado se repete.

Nunca foi observado um \dest{monopolo magnético} — um polo isolado. É uma
diferença importante em relação à eletricidade, onde a carga positiva pode
existir separada da negativa.
\atencao{Cortar um ímã não “separa” os polos: só cria mais ímãs. Com $n$
pedaços, ficam $n$ ímãs e $2n$ polos.}
\end{copiar}

\begin{figuranota}{Cada corte produz ímãs completos: de 1 ímã (2 polos) para 2
(4 polos) e para 4 (8 polos).}
\begin{tikzpicture}
  \pedaco{-3.0}{3.0}{1.15}
  \draw[naSuave, dashed, thin] (0,0.80) -- (0,1.50);
  \node[font=\tiny, color=naSuave, anchor=west] at (3.25,1.15) {1 ímã · 2 polos};

  \pedaco{-3.0}{-0.15}{0.15}  \pedaco{0.15}{3.0}{0.15}
  \node[font=\tiny, color=naSuave, anchor=west] at (3.25,0.15) {2 ímãs · 4 polos};

  \pedaco{-3.0}{-1.65}{-0.85} \pedaco{-1.35}{-0.15}{-0.85}
  \pedaco{0.15}{1.35}{-0.85}  \pedaco{1.65}{3.0}{-0.85}
  \node[font=\tiny, color=naSuave, anchor=west] at (3.25,-0.85) {4 ímãs · 8 polos};
\end{tikzpicture}
\end{figuranota}

% ============================================================
\section{Campo magnético e linhas de campo}

\begin{copiar}{Campo magnético $\vec{B}$}
\definicao O ímã modifica o espaço à sua volta: qualquer outro ímã colocado ali
sente uma força. Essa alteração do espaço é o \dest{campo magnético},
representado pelo vetor $\vec{B}$.

No SI, a unidade de $B$ é o \dest{tesla} (T). Também se usa o \dest{gauss} (G):
\[ 1\un{T} = 10^{4}\un{G} \]
\simbolos{$B$ (T); $1\un{G} = 10^{-4}\un{T}$.}
\diaadia{O campo da Terra vale cerca de $\dec{0,5}\un{G}$, ou
$5\times 10^{-5}\un{T}$ — um ímã de geladeira chega a $10^{-2}\un{T}$.}
\end{copiar}

\begin{copiar}{Linhas de campo}
\rotulo{Como se desenha.} As \dest{linhas de campo} indicam, em cada ponto, a
direção e o sentido de $\vec{B}$. Elas obedecem a três regras:
\begin{itens}
\item \dest{fora} do ímã, saem do polo \dest{norte} e entram no polo
      \dest{sul};
\item \dest{dentro} do ímã, vão do sul para o norte — as linhas são
      \dest{fechadas};
\item \dest{nunca se cruzam}, e onde estão mais \dest{próximas} o campo é mais
      \dest{intenso}.
\end{itens}
\atencao{Duas linhas cruzadas dariam duas direções para $\vec{B}$ no mesmo
ponto — impossível.}
\end{copiar}

\begin{figuranota}{Linhas de campo de um ímã em barra. Fora, vão de N para S;
perto dos polos elas se apertam, e ali o campo é mais forte.}
\begin{tikzpicture}[scale=0.80]
  \pedaco{-1.0}{1.0}{0}
  \foreach \h/\c in {0.95/1.75, 1.75/2.85, 2.60/4.05} {
    \draw[linhaB] (-1.0,0.25) .. controls (-\c,\h) and (\c,\h) .. (1.0,0.25);
    \draw[linhaB] (-1.0,-0.25) .. controls (-\c,-\h) and (\c,-\h) .. (1.0,-0.25);
  }
  \node[font=\tiny, color=naSuave] at (0,1.62) {$\vec{B}$};
\end{tikzpicture}
\end{figuranota}

% ============================================================
\section{A bússola e o campo da Terra}

\begin{copiar}{A Terra é um grande ímã}
\definicao A agulha da \dest{bússola} é um pequeno ímã que gira livremente:
ela se alinha com o campo magnético do lugar. Como a Terra tem campo próprio,
a agulha se orienta aproximadamente na direção norte–sul.

\rotulo{A troca de nomes.} O polo \dest{norte} da agulha é atraído pelo norte
geográfico. Como polos \dest{diferentes} se atraem, perto do norte geográfico
existe, na verdade, um polo \dest{sul magnético} da Terra.
\atencao{Norte geográfico $\neq$ norte magnético. O eixo magnético é inclinado
cerca de $11^\circ$ em relação ao eixo de rotação; o ângulo entre as duas
direções, num dado lugar, chama-se \dest{declinação magnética}.}
\end{copiar}

\begin{figuranota}{A Terra se comporta como um ímã de eixo inclinado. Junto ao
norte geográfico fica um polo sul magnético — por isso o N da agulha aponta
para lá.}
\begin{tikzpicture}[scale=0.90]
  \draw[naSuave, thick] (0,0) circle (1.15);
  % eixo geográfico
  \draw[naSuave, dashed, thin] (0,-1.55) -- (0,1.55);
  \node[font=\tiny, color=naSuave, anchor=south] at (0,1.58) {N geográfico};
  % ímã interno, inclinado
  \begin{scope}[rotate=-16]
    \fill[naAzul]  (-0.19,0) rectangle (0.19,0.92);
    \fill[naCoral] (-0.19,-0.92) rectangle (0.19,0);
    \node[font=\tiny\bfseries, text=naFundo] at (0,0.55)  {S};
    \node[font=\tiny\bfseries, text=naFundo] at (0,-0.55) {N};
  \end{scope}
  \node[font=\tiny, color=naAzul, anchor=west]  at (0.46,1.06) {S magnético};
  \node[font=\tiny, color=naCoral, anchor=east] at (-0.46,-1.06) {N magnético};
  % bússola
  \begin{scope}[shift={(2.55,0.20)}]
    \draw[naSuave, thin] (0,0) circle (0.42);
    \fill[naCoral] (-0.05,-0.02) -- (0.05,-0.02) -- (0,0.34) -- cycle;
    \fill[naAzul]  (-0.05,0.02) -- (0.05,0.02) -- (0,-0.34) -- cycle;
    \node[font=\tiny, color=naSuave, anchor=south] at (0,0.48) {bússola};
  \end{scope}
\end{tikzpicture}
\end{figuranota}

% ============================================================
\section{Materiais e imantação}

\begin{copiar}{Quem o ímã atrai}
\rotulo{Três comportamentos.}
\begin{itens}
\item \dest{Ferromagnéticos} — ferro, níquel, cobalto e suas ligas (aço):
      fortemente atraídos e capazes de \dest{ficar} imantados;
\item \dest{paramagnéticos} — alumínio, platina: atração fraquíssima;
\item \dest{diamagnéticos} — cobre, água, ouro: repulsão fraquíssima.
\end{itens}
Na prática, só os ferromagnéticos respondem de forma perceptível.
\end{copiar}

\begin{copiar}{Perder a imantação}
\rotulo{Como um ímã se estraga.} Aquecer ou golpear o ímã bagunça a
organização interna e ele perde a imantação. Acima de uma temperatura própria
de cada material — a \dest{temperatura de Curie} — o ferromagnetismo
desaparece.
\diaadia{Para o ferro, a temperatura de Curie é de cerca de $770\un{^\circ C}$:
um prego ao rubro deixa de ser atraído.}
\end{copiar}

\begin{exemplo}
O campo magnético da Terra, numa certa cidade, vale $\dec{0,50}\un{G}$.
\begin{alternativas}
\item Escreva esse valor em tesla.
\item Um ímã em barra é cortado ao meio, e cada metade é cortada de novo.
      Quantos ímãs e quantos polos existem no final?
\item Perto do polo norte geográfico, para onde aponta o polo norte da agulha
      de uma bússola? Por quê?
\end{alternativas}
\resolucao (a) Como $1\un{T} = 10^{4}\un{G}$, cada gauss vale $10^{-4}\un{T}$:
\[ B = \dec{0,50}\times 10^{-4}\un{T} \]
ou seja, \resultado{$B = \dec{5,0}\times 10^{-5}\un{T}$}.

(b) Dois cortes produzem \resultado{4 ímãs}, e cada um tem N e S:
\resultado{8 polos}. Nenhum pedaço fica com um polo só.

(c) Aponta \resultado{para o norte geográfico}, porque ali está um polo
\dest{sul} magnético da Terra — e polos diferentes se atraem.
\end{exemplo}

\begin{dica}
Três enganos frequentes: (i) achar que ímã atrai “todo metal” — alumínio e
cobre não são atraídos; (ii) trocar o norte geográfico pelo norte magnético;
(iii) desenhar as linhas de campo \dest{abertas}, parando no polo sul. Elas
continuam \dest{dentro} do ímã, de S para N, e se fecham.
\end{dica}

% ============================================================
\begin{exercicios}
\nivel{1}{Conceitos}
\begin{questoes}
\item Aproximando o polo norte de um ímã do polo norte de outro, observa-se
  \begin{alternativas}
  \item atração, pois todo ímã atrai qualquer outro.
  \item repulsão, pois polos iguais se repelem.
  \item nada, pois polos iguais não interagem.
  \item atração só se um deles for de ferro.
  \end{alternativas}
\item Partindo um ímã em barra exatamente ao meio, obtêm-se
  \begin{alternativas}
  \item um pedaço só com N e outro só com S.
  \item dois pedaços sem propriedades magnéticas.
  \item dois ímãs completos, cada um com N e S.
  \item um ímã e um pedaço de ferro comum.
  \end{alternativas}
\item A unidade de campo magnético no SI é
  \begin{alternativas}
  \item o gauss (G).
  \item o newton (N).
  \item o tesla (T).
  \item o ampère (A).
  \end{alternativas}
\item Fora de um ímã, as linhas de campo
  \begin{alternativas}
  \item saem do polo sul e entram no polo norte.
  \item saem do polo norte e entram no polo sul.
  \item partem dos dois polos e se cruzam no meio.
  \item formam linhas retas paralelas ao ímã.
  \end{alternativas}
\item Cite três materiais fortemente atraídos por um ímã e dois que não são.
\item Complete: $1\un{T} = \ldots\un{G}$, e o campo magnético da Terra vale
      cerca de $\ldots\un{T}$.
\end{questoes}

\nivel{2}{Aplicação}
\begin{questoes}
\item Um campo de $\dec{2,0}\un{G}$, em tesla, vale
  \begin{alternativas}
  \item $\dec{2,0}\times 10^{4}\un{T}$.
  \item $\dec{2,0}\times 10^{-4}\un{T}$.
  \item $\dec{2,0}\times 10^{-2}\un{T}$.
  \item $\dec{2,0}\times 10^{-6}\un{T}$.
  \end{alternativas}
\item A agulha da bússola aponta para o norte geográfico porque
  \begin{alternativas}
  \item lá existe um polo norte magnético da Terra.
  \item lá existe um polo sul magnético da Terra.
  \item a agulha é atraída pelo Sol.
  \item a Terra não tem campo magnético.
  \end{alternativas}
\item Um ímã é cortado em 3 pedaços iguais. O número de ímãs e de polos é
  \begin{alternativas}
  \item 3 ímãs e 3 polos.
  \item 3 ímãs e 6 polos.
  \item 1 ímã e 6 polos.
  \item 3 ímãs e 2 polos.
  \end{alternativas}
\item Passe $\dec{5,0}\times 10^{-5}\un{T}$ para gauss.
\item Numa região do desenho, as linhas de campo aparecem bem juntas. O que
      isso diz sobre $B$ ali? E se estivessem bem afastadas?
\item Um prego de ferro é aquecido até ficar ao rubro e aproximado de um ímã.
      O que se observa, e por quê?
\end{questoes}

\nivel{3}{Mais trabalhados}
\begin{questoes}
\item Um ímã em barra é cortado em $n$ pedaços. Escreva, em função de $n$, o
      número de ímãs e o número de polos obtidos. Existe algum $n$ que produza
      um polo isolado?
\item Numa cidade, $B = \dec{0,80}\un{G}$. Escreva esse valor em tesla e em
      microtesla ($1\un{\mu T} = 10^{-6}\un{T}$).
\item Sobre o monopolo magnético, é correto afirmar que
  \begin{alternativas}
  \item já foi obtido em laboratório, cortando ímãs muito finos.
  \item existe apenas em ímãs naturais, como a magnetita.
  \item nunca foi observado: todo pedaço de ímã tem os dois polos.
  \item é o nome do polo norte de um ímã em barra.
  \end{alternativas}
\item Explique, usando a definição de campo, por que duas linhas de campo não
      podem se cruzar.
\item Numa cidade a declinação magnética é de $21^\circ$ a oeste. Explique o
      que isso significa para quem se orienta por uma bússola e quer achar o
      norte geográfico.
\item Você recebe duas barras idênticas, sem marcação: uma é um ímã e a outra é
      ferro comum. Sem usar nenhum outro objeto, descreva um teste que
      identifique qual é o ímã, e justifique.
\end{questoes}

\gabarito{1b · 2c · 3c · 4b · 5) atraídos: ferro, níquel, cobalto (ou aço);
não atraídos: alumínio, cobre · 6) $10^{4}\un{G}$; $\approx 5\times
10^{-5}\un{T}$ · 7b · 8b · 9b · 10) $\dec{0,50}\un{G}$ · 11) juntas: campo mais
intenso; afastadas: mais fraco · 12) deixa de ser atraído — acima da
temperatura de Curie ($\approx 770\un{^\circ C}$) o ferro perde o
ferromagnetismo · 13) $n$ ímãs e $2n$ polos; nenhum $n$ dá polo isolado ·
14) $\dec{8,0}\times 10^{-5}\un{T} = 80\un{\mu T}$ · 15c · 16) o cruzamento
daria duas direções para $\vec{B}$ no mesmo ponto, e o campo tem uma só ·
17) o norte da bússola está $21^\circ$ a oeste do norte geográfico; para achar
o norte geográfico, corrija $21^\circ$ para leste · 18) encoste a ponta de uma
barra no \emph{meio} da outra e inverta os papéis: só há atração forte quando a
ponta usada é a do ímã, pois o meio do ímã é região neutra.}
\end{exercicios}

\end{document}
